% =====================================================================
\section{Testowanie rozwiązania}
\label{sec:testy}

PlayLink jest pokryty obszernym, zautomatyzowanym zestawem testów:
\textbf{135 testów} po stronie backendu (pytest) oraz \textbf{61 testów} po
stronie frontendu (Vitest). Pod względem \emph{różnorodności} wyróżniono
\textbf{sześć różnych rodzajów testów} (powyżej wymaganego minimum pięciu),
zestawionych w tab.~\ref{tab:typy-testow}. W podrozdziale~\ref{sec:katalog} opisano
ponadto \textbf{dziesięć reprezentatywnych przypadków testowych} (po kilka na rodzaj)
wraz z celem, metodą i wynikiem.

Pełna, rozszerzona dokumentacja testowa stanowi
\hyperref[sec:zalaczniki]{Załącznik B}. Obejmuje strategię i środowisko testowe,
22 szczegółowe przypadki, wyniki pokrycia i testów wydajnościowych oraz fragmenty
kodu testów pochodzące bezpośrednio z repozytorium.

\paragraph{Uwaga terminologiczna (różnorodność vs.\ liczba).}
Kryterium oceny dotyczy \emph{rodzajów} (typów) testów — wymaganych jest 5 różnych.
Liczba opisanych przypadków (10) to inna oś: są to konkretne testy rozłożone na owe
rodzaje (np.\ 2 na rodzaj). Oba wymagania są zatem spełnione jednocześnie: 6 rodzajów
oraz 10 szczegółowo udokumentowanych przypadków.

\renewcommand{\arraystretch}{1.3}
\begin{table}[H]
\centering\small
\begin{tabularx}{\textwidth}{@{}p{0.21\textwidth}p{0.18\textwidth}Yc@{}}
\toprule
\textbf{Rodzaj testów} & \textbf{Narzędzie} & \textbf{Zakres} & \textbf{Liczba} \\
\midrule
Jednostkowe & pytest, Vitest & Funkcje pomocnicze, walidacja, store'y, metryka Haversine & $\sim$25 \\
Integracyjne / API & FastAPI TestClient & Endpointy REST przez całą aplikację (SQLite in-memory) & $\sim$54 \\
WebSocket / realtime & TestClient WS & Rozgłaszanie ramek czatu i zdarzeń (34 połączenia WS) & $\sim$45 \\
Bezpieczeństwa & pytest + eth\_account & Replay, wygasły nonce, fałszywy admin, podrobione podpisy & $\sim$10 \\
Migracji (smoke) & pytest + Alembic & Migracja do \code{head}, weryfikacja schematu; w CI round-trip na Postgresie & 1 (+CI) \\
Kontraktowe (frontend) & Vitest & Server-loadery i akcje SvelteKit (\code{+page.server.ts}) & $\sim$20 \\
\bottomrule
\end{tabularx}
\caption{Sześć rodzajów testów w projekcie PlayLink.}
\label{tab:typy-testow}
\end{table}

\subsection{Charakterystyka rodzajów testów}

\paragraph{Testy jednostkowe.} Weryfikują izolowaną logikę bez bazy i sieci: parsowanie
listy adresów administratorów, funkcję \code{hash_nonce}, walidację nazw użytkownika,
a we frontendzie — podpisywanie, store'y oraz obliczanie odległości lobby (Haversine).
Narzędzia: pytest, Vitest + jsdom.

\paragraph{Testy integracyjne / API.} Uruchamiają żądania HTTP przez całą aplikację
FastAPI z użyciem \code{TestClient} i bazy SQLite in-memory (fixture nadpisuje
\code{get_session}). Pokrywają uwierzytelnianie, pełny cykl życia pokoju, panel
administratora i akcję kick — sprawdzając kody statusu oraz stan w bazie.

\paragraph{Testy WebSocket / realtime.} Wykorzystują \code{TestClient.websocket_connect}
(34 połączenia WS) do weryfikacji rozgłaszania ramek: \emph{event-update},
\emph{rsvp\_update}, \emph{roster\_update}, \emph{room\_closed}, \emph{member\_kicked}.

\paragraph{Testy bezpieczeństwa.} Najważniejszy rodzaj dla profilu projektu —
scenariusze ataków (replay, wygasły nonce, eskalacja uprawnień, podrobione podpisy
EIP-191), z realnym podpisywaniem kluczem (\code{eth_account}).

\paragraph{Testy migracji (smoke).} Aplikują pełny łańcuch Alembic
(\code{alembic upgrade head}) i sprawdzają obecność wszystkich ośmiu tabel; w CI
powtarzane na realnym PostgreSQL z round-tripem downgrade/upgrade i \code{alembic check}.

\paragraph{Testy kontraktowe frontendu.} Importują server-loadery i akcje SvelteKit
z zamockowanymi \code{fetch}/\code{cookies}, weryfikując przekierowania, strażników
autoryzacji i walidację formularzy.

\subsection{Katalog wybranych przypadków testowych}
\label{sec:katalog}

Poniżej opisano dziesięć reprezentatywnych przypadków testowych w jednolitym formacie
(rodzaj, cel, metoda/narzędzie, dane/kroki, oczekiwany wynik, rezultat). Wszystkie
przechodzą pomyślnie w pipeline CI; zespół może dołączyć do każdego zrzut ekranu
z wynikiem (np.\ raport pytest/Vitest lub log CI).

\newcommand{\tc}[7]{%
  \par\vspace{0.5em}\noindent\textbf{#1}\par\nopagebreak\vspace{0.2em}
  \noindent\begin{tabular}{@{}>{\itshape}p{0.20\textwidth}p{0.74\textwidth}@{}}
  Rodzaj & #2 \\
  Cel & #3 \\
  Metoda / narzędzie & #4 \\
  Dane / kroki & #5 \\
  Oczekiwany wynik & #6 \\
  Rezultat & #7 \\
  \end{tabular}\par\vspace{0.3em}\noindent\rule{\textwidth}{0.4pt}}

\tc{TC-01: Parsowanie adresów administratorów}
{Jednostkowy}
{Sprawdzić poprawne wczytanie i normalizację listy adresów z \code{ADMIN_ADDRESSES}.}
{pytest (\code{test_helpers.py}, \code{_parse_admin_addresses}).}
{Wejście \code{" 0xABC , 0xdef "} (spacje, różna wielkość liter).}
{Zbiór znormalizowanych adresów małymi literami: \code{{0xabc, 0xdef}}.}
{PASS}

\tc{TC-02: Hashowanie wartości nonce}
{Jednostkowy}
{Zapewnić, że nonce jest deterministycznie hashowany (SHA-256) przed zapisem.}
{pytest (\code{hash_nonce}).}
{Ten sam nonce hashowany dwukrotnie oraz dwa różne nonce.}
{Identyczny skrót dla identycznego wejścia; różne skróty dla różnych wejść.}
{PASS}

\tc{TC-03: Obliczanie odległości lobby (Haversine)}
{Jednostkowy (frontend)}
{Zweryfikować dobór pingu na podstawie odległości geograficznej lobby.}
{Vitest (\code{lobbyLocations.test.ts}).}
{Pary współrzędnych o znanej odległości.}
{Odległość w granicach tolerancji; właściwy przedział pingu.}
{PASS}

\tc{TC-04: Pełny przepływ logowania (challenge-response)}
{Integracyjny / API}
{Zweryfikować ścieżkę request-nonce $\rightarrow$ verify $\rightarrow$ wydanie JWT.}
{FastAPI \code{TestClient} (\code{test_auth.py}), podpis \code{eth_account}.}
{\code{POST /auth/request-nonce}, podpis nonce, \code{POST /auth/verify}.}
{HTTP 200 oraz poprawny token JWT i nazwa użytkownika w odpowiedzi.}
{PASS}

\tc{TC-05: Cykl życia pokoju i limit aktywnych pokojów}
{Integracyjny / API}
{Sprawdzić tworzenie, dołączanie, opuszczanie oraz limit 3 aktywnych pokojów.}
{FastAPI \code{TestClient} (\code{test_rooms_lifecycle.py}).}
{Sekwencja create/join/leave; próba utworzenia 4.\ pokoju.}
{Operacje 1--3 kończą się sukcesem; 4.\ pokój odrzucony błędem walidacji.}
{PASS}

\tc{TC-06: Rozgłaszanie wiadomości i ramek w czacie}
{WebSocket / realtime}
{Potwierdzić, że klienci otrzymują aktualizacje w czasie rzeczywistym.}
{\code{TestClient.websocket_connect} (\code{test_chat.py}, \code{test_room_events.py}).}
{Dwa połączenia WS w jednym pokoju; wysłanie wiadomości i zmiana RSVP.}
{Obaj klienci otrzymują ramki \emph{message}, \emph{rsvp-update}, \emph{roster-update}.}
{PASS}

\tc{TC-07: Odrzucenie ataku powtórzeniowego (replay)}
{Bezpieczeństwa}
{Uniemożliwić ponowne użycie tego samego nonce.}
{pytest (\code{test_verify_signature_rejects_replayed_nonce}).}
{Poprawna weryfikacja, a następnie ponowna z tym samym nonce/podpisem.}
{Pierwsza weryfikacja: 200; druga: odrzucenie (nonce oznaczony jako użyty).}
{PASS}

\tc{TC-08: Brak eskalacji uprawnień administratora}
{Bezpieczeństwa}
{Sfałszowane roszczenie o rolę admina nie może dawać dostępu.}
{pytest (\code{test_forged_admin_claim_does_not_grant_admin_access}).}
{Token/żądanie konta spoza \code{ADMIN_ADDRESSES} próbuje akcji administracyjnej.}
{HTTP 403; operacja administracyjna nie zostaje wykonana.}
{PASS}

\tc{TC-09: Migracja schematu bazy (smoke)}
{Migracji}
{Zapewnić, że pełny łańcuch migracji buduje poprawny schemat.}
{pytest + Alembic (\code{test_migrations.py}); w CI na PostgreSQL.}
{\code{alembic upgrade head} na świeżej bazie; inspekcja schematu.}
{Obecnych wszystkie 8 tabel i kluczowe kolumny; brak dryfu (\code{alembic check}).}
{PASS}

\tc{TC-10: Strażnik autoryzacji trasy (frontend)}
{Kontraktowy (frontend)}
{Niezalogowany użytkownik nie ma dostępu do tras chronionych.}
{Vitest (\code{rooms.name.page.server.test.ts}), zamockowane \code{cookies}/\code{fetch}.}
{Wywołanie \code{load} bez ważnej sesji.}
{Przekierowanie HTTP 303 do \code{/rooms}; brak ujawnienia danych chronionych.}
{PASS}

\subsection{Dokumentacja procesu i automatyzacja (CI)}
Proces testowy jest opisany w repozytorium (\code{docs/operations/testing.md} — tabele
pokrycia per plik, opis fixture'ów i kanałów WS). Konfiguracja: \code{pyproject.toml}
(\code{[tool.pytest.ini_options]}, pytest-cov), \code{vite.config.js} (Vitest, jsdom,
coverage-v8). Pipeline GitHub Actions (\code{.github/workflows/cicd.yml}) uruchamia
zadania: \code{backend-lint}, \code{backend-test} (z bramką \code{--cov-fail-under=80}),
\code{backend-migrations}, \code{frontend-lint}, \code{frontend-test} (z pokryciem) oraz
\code{build-test}. Zielony wynik wszystkich testów jest warunkiem wdrożenia
(\code{deploy}) oraz scalenia PR do \code{main}, co stanowi ciągłą, weryfikowalną
dokumentację wyników testowania.
